\section{Hyperbolic coordinates and the intrinsic boundary measure}
\label{sec:v14-hyperbolic}
The relative product law has a classical analytic model. We first
calculate that model in physical endpoint coordinates, and then identify
the boundary amplitudes intrinsically for the smooth billiards of
Theorem~\ref{thm:v8-main-relative}. These are complementary statements:
the first assumes normalizing charts; the second constructs a
one-dimensional linearizing coordinate directly from the already proved
half-line law. No smooth two-dimensional normal-form theorem is used in
the second argument.

\subsection{The analytic mixed-boundary calculation}
De Simoi--Kaloshin--Leguil~\cite{DSKL}, in the normal-form discussion on
printed page~10 of arXiv:1905.00890v4, recall the analytic symplectic
normal form near a hyperbolic period-two orbit. The local existence
statement is distinct from the symmetry assumptions in their geometric
inverse theorem; see also their discussion on printed page~13.
The following physical endpoint calculation was supplied in the
independent memorandum~\cite{A2ReviewV13NF}. We include its proof and
normalizations. The proposition is conditional on the indicated charts,
and is not attributed as a probability theorem to~\cite{DSKL}.

\begin{proposition}[Relative flux in analytic hyperbolic coordinates]
\label{prop:v14-normal-form}
Let
\begin{equation}\label{eq:v14-normal-map}
 N(s,p)=(\Delta(sp)s,\Delta(sp)^{-1}p),\qquad
 \Delta(0)=\lambda\in(0,1),
\end{equation}
where $\Delta$ is positive and analytic near zero. Let
\[
 \mathcal C_L(s,p)=(U(s,p),\mathcal P(s,p)),\qquad
 \mathcal C_R(q,t)=(V(q,t),\mathcal Q(q,t))
\]
be analytic canonical charts sending zero to zero, with
$U_s(0,0)V_t(0,0)\ne0$. For the local branch of
$\mathcal C_R\circ N^n\circ\mathcal C_L^{-1}$, prescribe the initial
physical position $u$ and terminal physical position $v$. There are a
fixed box in $(u,v)$, an integer $n_0$, and $\sigma\in(0,1)$ such that
this boundary problem has a unique local solution for all $n\ge n_0$.
Its physical mixed flux $J_n=|\partial P/\partial v|_u$ satisfies
\begin{equation}\label{eq:v14-normal-relative}
 \left\|\frac{J_n(u,v)}{J_n(0,0)}
                  -\mathcal B_L(u)\mathcal B_R(v)\right\|_{C^k}
                       \le C_k\sigma^n
\end{equation}
for every fixed $k$, where, writing
$U(\alpha(u),0)=u$ and $V(0,\beta(v))=v$,
\begin{equation}\label{eq:v14-projection-factors}
 \mathcal B_L(u)=\frac{U_s(0,0)}{U_s(\alpha(u),0)},\qquad
 \mathcal B_R(v)=\frac{V_t(0,0)}{V_t(0,\beta(v))}.
\end{equation}
These functions are positive and equal one at zero. If the physical
branch has a generating action $W_n$, normalized by
$P=-W_{n,u}$ and $Q=W_{n,v}$, then there are functions
$\mathcal S_L,\mathcal S_R$, vanishing with their first derivatives
at zero, for which
\begin{equation}\label{eq:v14-normal-action}
 \|W_n-W_n(0,0)-\mathcal S_L\oplus\mathcal S_R\|_{C^k}
                                                    \le C_k\sigma^n.
\end{equation}
All assertions hold with mixed parameter derivatives for a supplied
jointly smooth family of these functions and charts on common domains,
provided their fixed-order norms are uniformly bounded, their projection
derivatives are uniformly nonzero, and
$0<\lambda_-\le\lambda\le\lambda_+<1$.
\end{proposition}

\begin{proof}
Temporarily prescribe $s=s_0$ and $t=p_n$ rather than $u,v$.
The invariant $I=sp$ and the orbit endpoints obey
\begin{equation}\label{eq:v14-mixed-equation}
 I=st\Delta(I)^n,\qquad
 p_0=t\Delta(I)^n,\qquad q=s_n=s\Delta(I)^n.
\end{equation}
For a single map put $\lambda_+=\lambda$. Choose
$\lambda_+<\mu<\sigma<1$ and decrease the $I$ interval so
that $\Delta\le\mu$. On a fixed sufficiently small $(s,t)$ box,
$|st|n\mu^{n-1}\|\Delta'\|_\infty<1/2$ for every $n\ge1$,
since $\sup_n n\mu^{n-1}<\infty$. The right side of the first
identity in \eqref{eq:v14-mixed-equation} maps a common small closed
$I$ interval into itself. It is a contraction there. Its solution
satisfies $|I|\le C\mu^n$. The intermediate coordinates are
$(s\Delta(I)^i,t\Delta(I)^{n-i})$, so the whole orbit remains
inside the chosen local chart.

Set
\[
 a_n=\Delta(I)^n,\qquad
 D_n=1-nI\frac{\Delta'(I)}{\Delta(I)}.
\]
The contraction estimate gives $|D_n|\ge1/2$. Differentiating
\eqref{eq:v14-mixed-equation}, without dividing by either endpoint,
gives
\begin{equation}\label{eq:v14-exact-mixed-flux}
 T_n:=\left.\frac{\partial p_0}{\partial t}\right|_s
      =\left.\frac{\partial q}{\partial s}\right|_t
      =\frac{\Delta(I)^n}{D_n}.
\end{equation}
For example, $I_t=sa_n/D_n$ and
$(p_0)_t=a_n+tn\Delta(I)^{n-1}\Delta'(I)I_t=a_n/D_n$.
This also proves the formula on both axes.

Here are the differentiated relative estimates. At every fixed order,
derivatives of the function $\Delta(I)^n$ with $I$ treated as an
independent variable are bounded by $C_k(1+n)^k\mu^n$; a finite
number of small $n$ are absorbed in $C_k$. Implicit differentiation
of the first equation in \eqref{eq:v14-mixed-equation} leaves $D_n$
on the highest derivative of $I$. Its remaining terms contain the
preceding derivatives of $I$, derivatives of $st$, and derivatives of
$\Delta^n$. Induction therefore gives
\[
 \|I\|_{C^k}\le C_k(1+n)^{m_k}\mu^n.
\]
These norms can include the supplied parameters. Since
$\log(\Delta(I)/\lambda)=Ih(I)$ with a jointly smooth bounded
function $h$, the same induction gives
\[
 \left\|n\log\frac{\Delta(I)}\lambda\right\|_{C^k}
       +\|D_n-1\|_{C^k}\le C_k(1+n)^{r_k}\mu^n.
\]
The identity at $I=0$ holds throughout the parameter family; thus
parameter derivatives of $\lambda$ do not leave an unsuppressed
term in this normalized expression. Exponentiating and using
$|D_n|\ge1/2$ proves
\begin{equation}\label{eq:v14-normalized-T}
             \|\lambda^{-n}T_n-1\|_{C^k}\le C_k\sigma^n.
\end{equation}
One fixed strict margin $\mu<\sigma$ absorbs the polynomial at
each fixed order; the constants, not the margin, depend on that order.

The physical boundary map in the mixed coordinates is
\begin{equation}\label{eq:v14-Phi}
 \Phi_n(s,t)=(U(s,t a_n),V(s a_n,t)),\qquad
 \Phi_\infty(s,t)=(U(s,0),V(0,t)).
\end{equation}
It converges to $\Phi_\infty$ in every fixed derivative norm at an
exponential rate. Choose fixed inner and outer boxes on which the two
limiting scalar projections are diffeomorphisms, with the closure of
the inner box strictly inside the image of the outer one. Conjugating
$\Phi_n$ by $\Phi_\infty^{-1}$ gives the identity plus a uniformly
small $C^1$ perturbation. The equation for its inverse is a contraction
on a small ball about each point of the inner box, uniformly in that
point. This gives a unique inverse on a common nonshrinking physical
box for all $n\ge n_0$. Subtracting the inverse equations gives its
$C^0$ convergence. Differentiation isolates the same uniformly
invertible differential and proves the $C^k$ convergence by induction.
In particular $(s,t)=\Phi_n^{-1}(u,v)$ tends to
$(\alpha(u),\beta(v))$ at the stated rate.

Put $\mathcal D_n=\det D\Phi_n$. Canonicity of the initial chart
and \eqref{eq:v14-exact-mixed-flux} give, on the boundary surface,
\[
 du\wedge dP=ds\wedge dp_0=T_n\,ds\wedge dt,
 \qquad du\wedge dv=\mathcal D_n\,ds\wedge dt.
\]
Consequently the exact physical, rather than canonical, flux is
\begin{equation}\label{eq:v14-physical-flux}
 J_n(u,v)=\left|\frac{T_n(s,t)}{\mathcal D_n(s,t)}\right|,
                         \qquad (s,t)=\Phi_n^{-1}(u,v).
\end{equation}
At zero this denominator is exactly
\begin{equation}\label{eq:v14-reference-determinant}
 \mathcal D_n(0,0)=U_s(0,0)V_t(0,0)
              -\lambda^{2n}U_p(0,0)V_q(0,0).
\end{equation}
On the physical box its limit is
$U_s(\alpha(u),0)V_t(0,\beta(v))$, uniformly bounded away from
zero. Equations~\eqref{eq:v14-normalized-T}--\eqref{eq:v14-reference-determinant}
therefore prove \eqref{eq:v14-normal-relative} and
\eqref{eq:v14-projection-factors}, including all fixed mixed derivatives.
This step is necessary: canonical boundary coordinates are not physical
configuration observations.

Finally define
\[
 \mathcal S_L(u)=-\int_0^u\mathcal P(\alpha(x),0)\,\dd x,
 \qquad
 \mathcal S_R(v)=\int_0^v\mathcal Q(0,\beta(y))\,\dd y.
\]
The physical first variations of $W_n$ converge to the two displayed
integrands, with all fixed derivatives, by \eqref{eq:v14-Phi}.
Integrate the difference of the gradients along the two coordinate
segments from $(0,0)$ to $(u,v)$. The values at the origin fix the
additive constants and prove \eqref{eq:v14-normal-action}.
\end{proof}

For the two-step billiard return the contracting eigenvalue in
\eqref{eq:v14-normal-map} is $\lambda=e^{-2\gamma}$, and $n$
returns mean $j=2n$ flights. At contact $b$ its stable and unstable
linear directions in physical position--momentum coordinates are
$(1,-a_b)$ and $(1,a_b)$, by \eqref{eq:g-effective}. They are
nonvertical, regardless of equality or inequality of the curvatures.
Use the same canonical chart at both ends. Write its first derivative
as
\[
 D\mathcal C(0)=
 \begin{pmatrix}A&B\\-a_bA&a_bB\end{pmatrix}.
\]
Canonicity gives $AB=(2a_b)^{-1}$. Thus
\eqref{eq:v14-reference-determinant} and \eqref{eq:v14-physical-flux}
recover the exact physical normalization
\begin{equation}\label{eq:v14-csch-normalization}
 J_n(0,0)=\frac{2a_b\lambda^n}{1-\lambda^{2n}}
                        =\frac{a_b}{\sinh(2n\gamma)}.
\end{equation}
This is an exact agreement with the Jacobi formula, not merely an
agreement of leading exponential orders. If both constructions apply
to the same analytic billiard, uniqueness of the normalized limit and
evaluation at the other endpoint zero give
$B_b=\mathcal B_L=\mathcal B_R$. Their action limits agree as well,
with the stated zero normalizations. The normal-form function $\Delta$
alone does not specify the physical amplitudes: the projection charts
enter \eqref{eq:v14-projection-factors} essentially.

\subsection{A smooth intrinsic characterization}
We now return to the full smooth setting. Write $W_{j,b}$ for the
stationary action starting at type $b$, and let
\[
 \varphi_b(u)=x_1^{(b)}(u),\qquad
 r_b=\sqrt{\frac{c_b}{c_{1-b}}}\,e^{-\gamma},\qquad
 R_b=\varphi_{1-b}\circ\varphi_b,\qquad
 \lambda=e^{-2\gamma}.
\]
Here $x^{(b)}$ is the physical half-line in
Lemma~\ref{lem:v4-halfline}. Thus $\varphi_b$ is the first-hit
position along the stable branch, not a freely chosen change of
variables. Notice that $r_0r_1=\lambda$ and $0<r_b<1$.
Indeed $\sqrt{c_b/c_{1-b}}e^{-\gamma}
=c_b/(c_0c_1+\sqrt{c_0c_1(c_0c_1-1)})<1$, since $c_{1-b}>1$.

\begin{theorem}[The boundary amplitude as a linearizing density]
\label{thm:v14-intrinsic-density}
On a fixed smaller collar, uniformly for the smooth compact families
of Theorem~\ref{thm:v8-main-relative}, the half-line amplitudes satisfy
\begin{equation}\label{eq:v14-one-flight-cocycle}
 B_{1-b}(\varphi_b(u))\varphi_b'(u)=r_bB_b(u),
                              \qquad b=0,1.
\end{equation}
The increasing smooth coordinates
\begin{equation}\label{eq:v14-linearizing-coordinate}
 \zeta_b(u)=\int_0^u B_b(x)\,\dd x
\end{equation}
are normalized by $\zeta_b(0)=0$, $\zeta_b'(0)=1$ and obey
\begin{equation}\label{eq:v14-linearization}
 \zeta_{1-b}(\varphi_b(u))=r_b\zeta_b(u),\qquad
 \zeta_b(R_b(u))=\lambda\zeta_b(u).
\end{equation}
For each $b$, $\zeta_b$ is the unique differentiable local coordinate
with that normalization which linearizes $R_b$. Moreover, for a fixed
$\sigma\in(\lambda_+,1)$, where $\lambda_+$ is the maximum of
$\lambda$ on the compact family,
\begin{equation}\label{eq:v14-iterate-limits}
 \|\lambda^{-n}R_b^n-\zeta_b\|_{C^k}
 +\|\lambda^{-n}(R_b^n)'-B_b\|_{C^k}
                                  \le C_k\sigma^n.
\end{equation}
The norms include every fixed mixed parameter and endpoint derivative.
No normalizing chart for the two-dimensional return map is a hypothesis.
\end{theorem}

\begin{proof}
Put $o=1-b$. Uniqueness of the half-line tail gives the exact relations
\begin{equation}\label{eq:v14-Bellman}
 S_b(u)=\ell_b(u,\varphi_b(u))-g+S_o(\varphi_b(u)),\qquad
 \ell_{b,2}(u,\varphi_b(u))+S_o'(\varphi_b(u))=0.
\end{equation}
The second relation also follows from first variation of the middle
endpoint. Its $w$ derivative is
$\ell_{b,22}(u,w)+S_o''(w)$, uniformly positive on the small box.
The implicit function theorem and the already constructed half-line
therefore give
\begin{equation}\label{eq:v14-first-hit-derivative}
 \varphi_b'(u)=
 \frac{-\ell_{b,12}(u,\varphi_b(u))}
 {\ell_{b,22}(u,\varphi_b(u))+S_o''(\varphi_b(u))}>0.
\end{equation}
At zero this is
$(c_o+g a_o)^{-1}=\sqrt{c_b/c_o}e^{-\gamma}=r_b$.
After decreasing the common interval, both first-hit maps are
contractions into it. All their mixed derivatives are controlled by
the same positive denominator. The graph of stable initial physical
states is $P=-S_b'(u)$: differentiating the first identity in
\eqref{eq:v14-Bellman} gives $P=-\ell_{b,1}$, while the terminal
momentum is $\ell_{b,2}=-S_o'$. The reversed unstable graph is
$P=S_b'(u)$. This verifies the dynamical interpretation directly
from the physical generating action.

We prove the amplitude identity from exact finite concatenation.
For a tail of $j$ flights starting at $o$, let $w_j=w_j(u,v)$ solve
\[
 \ell_{b,2}(u,w_j)+W_{j,o,1}(w_j,v)=0.
\]
Strict convexity identifies the resulting stationary value with
$W_{j+1,b}(u,v)$. Set
$H_j=\ell_{b,22}(u,w_j)+W_{j,o,11}(w_j,v)$.
Eliminating the middle endpoint gives the exact Schur identity
\begin{equation}\label{eq:v14-concatenation}
 -W_{j+1,b,12}(u,v)=
 \frac{-\ell_{b,12}(u,w_j)}{H_j}
                          [-W_{j,o,12}(w_j,v)].
\end{equation}
The denominator is uniformly positive. Theorem~\ref{thm:v4-factorization}
and its differentiated action limit imply
$w_j\to\varphi_b(u)$ and
$H_j\to\ell_{b,22}(u,\varphi_b(u))+S_o''(\varphi_b(u))$
on the fixed box. The inverse equations, with the same positive
denominator, give this convergence in every fixed mixed norm.

The terminal type of both sides of \eqref{eq:v14-concatenation}
is $c=(o+j)\bmod2$. The exact reference constants are
\[
 d^0_{j,o}=\frac{\sqrt{a_oa_c}}{\sinh(j\gamma)},\qquad
 d^0_{j+1,b}=\frac{\sqrt{a_ba_c}}{\sinh((j+1)\gamma)}.
\]
Their ratio tends to
$\sqrt{a_o/a_b}e^\gamma=r_b^{-1}$, with all fixed parameter
derivatives. Divide \eqref{eq:v14-concatenation} by
$d^0_{j+1,b}$ and let $j$ tend to infinity along either parity.
Relative factorization and \eqref{eq:v14-first-hit-derivative} give
\[
 B_b(u)B_c(v)=r_b^{-1}\varphi_b'(u)
                   B_o(\varphi_b(u))B_c(v).
\]
Since $B_c>0$, cancellation proves
\eqref{eq:v14-one-flight-cocycle}. Integration from zero proves
\eqref{eq:v14-linearization}. This argument uses the previously
proved relative theorem to identify its amplitude; it does not claim
that a scalar linearization alone proves the finite-bridge theorem.

For uniqueness, conjugate any other normalized differentiable
linearizer by $\zeta_b^{-1}$. The resulting function $h$ has
$h(0)=0$, $h'(0)=1$ and $h(\lambda z)=\lambda h(z)$. Hence
$h(z)=\lambda^{-n}h(\lambda^nz)\to z$, so $h(z)=z$.
For the quantitative assertion, the exact identity is
\[
 R_b^n(u)=\zeta_b^{-1}(\lambda^n\zeta_b(u)).
\]
Write $\zeta_b^{-1}(x)=x+x^2h_b(x)$, where $h_b$ is jointly
smooth on a common interval. Division by $\lambda^n$ leaves
$\zeta_b(u)+\lambda^n\zeta_b(u)^2h_b(\lambda^n\zeta_b(u))$.
At each fixed derivative order this last term, including parameter
derivatives of $\lambda$, is bounded by
$C_k(1+n)^{m_k}\lambda_+^n$. Its endpoint derivative proves the
second estimate in \eqref{eq:v14-iterate-limits}; the fixed strict
margin gives the displayed rate. This also proves all claimed
uniformities without extending the scalar coordinate to a symplectic
chart off the invariant branch.
\end{proof}

\begin{corollary}[Stable-action width determined by the even laws]
\label{cor:v14-width}
Put $\widehat S_b=S_b\circ\zeta_b^{-1}$. Then the limiting laws
can be written
\begin{equation}\label{eq:v14-flat-measure-law}
 \mathcal F_{bc}(d)=\frac{\sqrt{a_ba_c}}{\pi d^2}
       \int(d-\widehat S_b(x)-\widehat S_c(y))_+\,\dd x\dd y.
\end{equation}
For small $E>0$, let $u_{b,-}(E)<0<u_{b,+}(E)$ solve $S_b(u)=E$.
The two even laws determine the normalized widths
\begin{equation}\label{eq:v14-width}
 \begin{split}
 \mathcal W_b(E)
   &:=\sqrt{\frac{a_b}{2}}
       [\zeta_b(u_{b,+}(E))-\zeta_b(u_{b,-}(E))]\\
   &=\int_0^E x^{-1/2}\mathcal V_b(x)\,\dd x.
 \end{split}
\end{equation}
Thus the smooth invariant is the complete normalized sublevel width
of the half-line action in its intrinsic linearizing coordinate.
\end{corollary}
\begin{proof}
The change of variables $x=\zeta_b(u)$ has differential
$\dd x=B_b(u)\dd u$. This proves \eqref{eq:v14-flat-measure-law}.
The pushforward calculation in \eqref{eq:v9-energy-profile} gives
\[
 \int_{S_b(u)<E}B_b(u)\,\dd u
       =\sqrt{\frac2{a_b}}\int_0^E x^{-1/2}\mathcal V_b(x)\,\dd x.
\]
It is also exactly
$\zeta_b(u_{b,+}(E))-\zeta_b(u_{b,-}(E))$.
Theorem~\ref{thm:v9-compatibility} determines $\mathcal V_b$ from
the even law, proving the assertion. That theorem's proof uses only
the half-line formula and Volterra uniqueness, not this corollary.
\end{proof}

The normalization in \eqref{eq:v14-width} is useful when the separate
curvatures are not supplied. The even laws determine $\mathcal W_b$
without knowing $a_b$ separately. With $a_b$ supplied, they also
determine the unnormalized coordinate width. When both participating contact graphs are even,
$S_b$ and $B_b$ are even and $\zeta_b$ is odd, so this width specifies
both inverse branches of $\widehat S_b$. For a general contact it
specifies their difference. The physical contact-graph inverse in
Section~\ref{sec:v12-two-contact} is a further geometric argument,
not an inference that a symmetrized width determines an arbitrary
asymmetric boundary.

\subsection{The comparison at the level of the physical theorem}
Proposition~\ref{prop:v14-normal-form} explains the analytic relative
product mechanism, including its physical projection and exact
reference normalization. If the limiting endpoint Hessians are
positive, Proposition~\ref{prop:v9-morse-transport} also transports
its estimates through the normalized residual integral. At an isolated
analytic billiard channel the positivity and first-hit property are
verified in Sections~\ref{sec:g-localization}--\ref{sec:g-action}.
An odd-flight analytic comparison can likewise compose the terminal
chart with one physical flight. Thus neither the parity distinction nor
the absence of global finite horizon is an obstruction to that local
analytic mechanism. They are requirements handled simultaneously in
our physical theorem, not separate claims of a new normal-form effect.

Theorem~\ref{thm:v8-main-relative} does not supply normalizing charts
as data or hypotheses. It starts with smooth obstacle graphs and
positive geometric margins. Its half-line and finite-bridge
contractions, localized trace-ideal estimates, and common sublevel
charts prove the mixed parameter and offset bounds on one closed
collar, with the exact phase normalization and the physical event
selection. A pointwise analytic normal form supplies none of these
uniform chart estimates for an arbitrary smooth family. The conditional
family clause in Proposition~\ref{prop:v14-normal-form} states exactly
what would be needed to use that route instead. No assertion that such
a smooth normal-form route is impossible, or absent from all prior
literature, is required here.

Theorem~\ref{thm:v14-intrinsic-density} makes the geometric meaning of
the determinant amplitude explicit in the full smooth class: it is
the derivative of the normalized scalar linearizer along the physical
stable branch, with an exact first-flight cocycle linking the two
contact types. Corollary~\ref{cor:v14-width} identifies what the
complete invariant measures in that coordinate. These conclusions
connect the analytic interpretation to the constructed smooth
amplitude rather than counting the interpretation as an unrelated
rigidity theorem.

The remaining steps concern information that the forward benchmark
does not reconstruct: the Volterra inverse for the complete widths;
the explicit independent-contact geometric inverse, its analytic
physical image and its two-flight coordinates; and the regularized
binary experiment with its charged smooth-class calibration.
The preparation theorem uses certified physical smoothness bounds,
not access to $\Delta$ or to a canonical chart. Its sufficient rate
is not claimed to be optimal. The central result is the joint passage
from smooth physical records to an intrinsic boundary measure, to its
complete observable invariant, and to a finite-preparation inverse.
The classical analytic product mechanism is part of that explanation,
not the novelty claimed for the package.
